% GENERATED FILE. Edit canonical lessons, then run npm run generate:books.
% canonical-source-hash: fnv1a64:469ee4ea64329466
% canonical-lessons: ZH-C403-houcheshi, ZH-C403-luke, ZH-W403-lu-travel, ZH-C403-zhengdian, ZH-C403-liecheyuan, ZH-C403-zhanzhang

\chapter{Travellers}
\label{ch:zh-travellers}
\hlchaptermodality{403}

\begin{chapteropening}
\textbf{By the end of this chapter:} \emph{I can say we were in time for the train, talk about travellers and trains on time, and find the staff.}

Complete the last lesson of chapter 403: I can say we were in time for the train, talk about travellers and trains on time, and find the staff.
\end{chapteropening}

\section[hòuchēshì]{hòuchēshì — station waiting room}

\label{lesson:ZH-C403-houcheshi}



\begin{quote}
Before the new one: say the Mandarin for to change trains, then the Mandarin for a railway carriage.
\end{quote}

\subsection*{You'll want to know: hòuchēshì}
\emph{hòuchēshì} — \textquotedblleft{}station waiting room\textquotedblright{}.

\begin{quote}
\emph{nǐ zuò zài hòuchēshì děng wǒ} — Sit in the waiting room and wait for me.
\end{quote}

Its characters come later; for now it is a word for the ear.

\begin{cousinweb}
\emph{hòuchēshì} is built from \zh{候车} 'to wait for a train' and \emph{shì} 'a room', all inherited from Middle Chinese, \emph{huwH}, \emph{tsyhae} and \emph{syit} (Guangyun, in Baxter's transcription), which glosses the last 'a room'.
\end{cousinweb}

\begin{grammarlens}[title={What you've built}]
One more word for getting around.
\end{grammarlens}

\subsection*{Your turn}
\begin{itemize}
\raggedright
  \item \emph{Say it:} \emph{hòuchēshì}
  \item \emph{Say it:} \emph{hòuchēshì}, once more
  \item \emph{Say it:} \emph{nǐ zuò zài hòuchēshì děng wǒ} — Sit in the waiting room and wait for me.
  \item \emph{Recall:} say the Mandarin for to change trains, then the Mandarin for a railway carriage, then say \emph{hòuchēshì} again
\end{itemize}

\begin{tcolorbox}[breakable,colback=teal!4,colframe=teal!35!black,title={Before you move on}]
What does hòuchēshì mean? (\textquotedblleft{}Station waiting room\textquotedblright{}.) Say it once more.
\end{tcolorbox}
\section[\texorpdfstring{\zh{旅客} (lǚkè)}{旅客 (lǚkè)}]{\zh{旅客} (lǚkè) — traveller; passenger on a long journey}

\label{lesson:ZH-C403-luke}



\begin{quote}
Before the new one: say the Mandarin for a railway carriage, then the Mandarin for a station waiting room.
\end{quote}

\subsection*{You'll want to know: \zh{旅客}}
\textbf{\zh{旅客}} — \emph{lǚkè} — \textquotedblleft{}traveller; passenger on a long journey\textquotedblright{}.

\begin{quote}
\textbf{\zh{旅客都上车了}} \emph{lǚkè dōu shàngchē le} — The passengers have all got on.
\end{quote}

The new character, in pieces: \zh{旅} — ten strokes in two pieces: on the left a dot, a bar, a stroke that runs right and turns down with a hook, and a long stroke falling to the left; on the right a short stroke falling to the left and a bar, then a short stroke falling to the left, a vertical that turns up at the foot, a short stroke falling to the left and a long stroke falling to the right. The next lesson puts it in your hand.

\begin{cousinweb}
\textbf{\zh{旅客}} (\emph{lǚkè}) is built from \zh{旅} and \zh{客} 'a guest', both inherited from Middle Chinese, \emph{ljoX} and \emph{khaek} (Guangyun, in Baxter's transcription). The dictionary glosses the first as a body of troops, five hundred men after an older book; the sense 'to travel' is not given there, so how it came is not settled.
\end{cousinweb}

\begin{grammarlens}[title={What you've built}]
One more word for getting around.
\end{grammarlens}

\subsection*{Your turn}
\begin{itemize}
\raggedright
  \item \emph{Say it:} \emph{lǚkè}
  \item \emph{Say it:} \emph{lǚkè}, once more
  \item \emph{Say it:} \emph{lǚkè dōu shàngchē le} — The passengers have all got on.
  \item \emph{Recall:} say the Mandarin for a railway carriage, then the Mandarin for a station waiting room, then say \emph{lǚkè} again
\end{itemize}

\begin{tcolorbox}[breakable,colback=teal!4,colframe=teal!35!black,title={Before you move on}]
What does \zh{旅客} mean? (\textquotedblleft{}Traveller; passenger on a long journey\textquotedblright{}.) Say it once more.
\end{tcolorbox}
\section[\texorpdfstring{\zh{旅} (lǚ)}{旅 (lǚ)}]{\zh{旅} — a left piece of four strokes, then a right piece of six}

\label{lesson:ZH-W403-lu-travel}



\begin{quote}
The word first, then the character inside it.

\begin{itemize}
\raggedright
  \item \emph{Recall:} say \emph{lǚkè} — \textbf{R1}, one lesson back, before it decays
\end{itemize}

\begin{quote}
\zh{客} \zh{走}
\end{quote}

\emph{Read it:} each one aloud, with nothing said first — \emph{kè}, \emph{zǒu}
\end{quote}

\subsection*{Script — a left piece of four strokes, then a right piece of six}
\hlblockfigure{\detokenize{figures/ZH-W403-lu-travel-filmstrip.pdf}}{How \zh{旅} is written, stroke by stroke}

\begin{quote}
\zh{旅}
\end{quote}

\textbf{Two pieces, left before right.} Count the pieces, not the strokes.

\textbf{The left piece, four strokes.} A dot on top. A bar under it. A stroke that runs right, turns down and hooks to the left at the foot. Then a long stroke falling to the left from under the bar.

\textbf{The right piece, six strokes.} A short stroke falling to the left and a bar to its right, like a small roof. Under them a short stroke falling to the left, a vertical that turns up to the right at the foot, another short stroke falling to the left, and last a long stroke falling to the right.

4 + 6 = ten.

\subsection*{Your turn}
\begin{itemize}
\raggedright
  \item \emph{Write it:} \zh{旅}, all ten, with the model in front of you
  \item \emph{Cover:} write it again from memory, and count the pieces, not the strokes
\end{itemize}

\begin{tcolorbox}[breakable,colback=teal!4,colframe=teal!35!black,title={Before you move on}]
How many pieces does \zh{旅} have? (\textbf{Two: four strokes and six}.) What is the last stroke? (\textbf{A long stroke falling to the right}.)

Source: \href{https://raw.githubusercontent.com/chanind/hanzi-writer-data/68d10a4b21150cae5e1ebbd223eed289cf32d90c/data/\%E6\%97\%85.json}{Hanzi Writer Data \zh{旅}, pinned stroke paths and medians}.
\end{tcolorbox}
\section[\texorpdfstring{\zh{正点} (zhèngdiǎn)}{正点 (zhèngdiǎn)}]{\zh{正点} (zhèngdiǎn) — on schedule (of a train or plane)}

\label{lesson:ZH-C403-zhengdian}



\begin{quote}
Before the new one: say the Mandarin for a station waiting room, then the Mandarin for a passenger on a long journey.
\end{quote}

\subsection*{You'll want to know: \zh{正点}}
\textbf{\zh{正点}} — \emph{zhèngdiǎn} — \textquotedblleft{}on schedule (of a train or plane)\textquotedblright{}.

\begin{quote}
\textbf{\zh{我们的火车正点到了}} \emph{wǒmen de huǒchē zhèngdiǎn dào le} — Our train got in on time.
\end{quote}

\begin{cousinweb}
\textbf{\zh{正点}} (\emph{zhèngdiǎn}) is built from \zh{正} 'exactly' and \zh{点} 'o'clock', both inherited from Middle Chinese, \emph{tsyengH} and \emph{temX} (Guangyun, in Baxter's transcription), which glosses the first 'right, proper' and the second as a dot made with the brush: on the hour the timetable gives.
\end{cousinweb}

\begin{grammarlens}[title={What you've built}]
One more word for getting around.
\end{grammarlens}

\subsection*{Your turn}
\begin{itemize}
\raggedright
  \item \emph{Say it:} \emph{zhèngdiǎn}
  \item \emph{Say it:} \emph{zhèngdiǎn}, once more
  \item \emph{Say it:} \emph{wǒmen de huǒchē zhèngdiǎn dào le} — Our train got in on time.
  \item \emph{Recall:} say the Mandarin for a station waiting room, then the Mandarin for a passenger on a long journey, then say \emph{zhèngdiǎn} again
\end{itemize}

\begin{tcolorbox}[breakable,colback=teal!4,colframe=teal!35!black,title={Before you move on}]
What does \zh{正点} mean? (\textquotedblleft{}On schedule (of a train or plane)\textquotedblright{}.) Say it once more.
\end{tcolorbox}
\section[lièchēyuán]{lièchēyuán — train attendant}

\label{lesson:ZH-C403-liecheyuan}



\begin{quote}
Before the new one: say the Mandarin for a passenger on a long journey, then the Mandarin for on schedule.
\end{quote}

\subsection*{You'll want to know: lièchēyuán}
\emph{lièchēyuán} — \textquotedblleft{}train attendant\textquotedblright{}.

\begin{quote}
\emph{zhǎo bu dào zuòwèi dehuà kěyǐ wèn lièchēyuán} — If you can't find your seat, ask the train attendant.
\end{quote}

Its characters come later; for now it is a word for the ear.

\begin{cousinweb}
\emph{lièchēyuán} is built from \emph{liè} 'a row', \zh{车} 'vehicle' and \emph{yuán} 'a member of staff', all inherited from Middle Chinese, \emph{ljet}, \emph{tsyhae} and \emph{hjwen} (Guangyun, in Baxter's transcription). The dictionary glosses the first 'a row, an order' and the third as the number of things; a train is a row of carriages, and the sense 'staff' for the third is not given there.
\end{cousinweb}

\begin{grammarlens}[title={What you've built}]
One more word for getting around.
\end{grammarlens}

\subsection*{Your turn}
\begin{itemize}
\raggedright
  \item \emph{Say it:} \emph{lièchēyuán}
  \item \emph{Say it:} \emph{lièchēyuán}, once more
  \item \emph{Say it:} \emph{zhǎo bu dào zuòwèi dehuà kěyǐ wèn lièchēyuán} — If you can't find your seat, ask the train attendant.
  \item \emph{Recall:} say the Mandarin for a passenger on a long journey, then the Mandarin for on schedule, then say \emph{lièchēyuán} again
\end{itemize}

\begin{tcolorbox}[breakable,colback=teal!4,colframe=teal!35!black,title={Before you move on}]
What does lièchēyuán mean? (\textquotedblleft{}Train attendant\textquotedblright{}.) Say it once more.
\end{tcolorbox}
\section[\texorpdfstring{\zh{站长} (zhànzhǎng)}{站长 (zhànzhǎng)}]{\zh{站长} (zhànzhǎng) — stationmaster; head of a station}

\label{lesson:ZH-C403-zhanzhang}



\begin{quote}
Before the new one: say the Mandarin for on schedule, then the Mandarin for a train attendant.
\end{quote}

\subsection*{You'll want to know: \zh{站长}}
\textbf{\zh{站长}} — \emph{zhànzhǎng} — \textquotedblleft{}stationmaster; head of a station\textquotedblright{}.

\begin{quote}
\textbf{\zh{旅客去找站长}} \emph{lǚkè qù zhǎo zhànzhǎng} — The passenger goes to find the stationmaster.
\end{quote}

\begin{cousinweb}
\textbf{\zh{站长}} (\emph{zhànzhǎng}) is built from \zh{站} 'a station' and \zh{长} 'the head', both inherited from Middle Chinese, \emph{treamH} and \emph{trjangX} (Guangyun, in Baxter's transcription). The dictionary glosses the first as an everyday word for standing alone and the second 'big'; a station is usually explained as a place where vehicles stand.
\end{cousinweb}

\begin{grammarlens}[title={What you've built}]
One more word for getting around.
\end{grammarlens}

\subsection*{Your turn}
\begin{itemize}
\raggedright
  \item \emph{Say it:} \emph{zhànzhǎng}
  \item \emph{Say it:} \emph{zhànzhǎng}, once more
  \item \emph{Say it:} \emph{lǚkè qù zhǎo zhànzhǎng} — The passenger goes to find the stationmaster.
  \item \emph{Recall:} say the Mandarin for on schedule, then the Mandarin for a train attendant, then say \emph{zhànzhǎng} again
\end{itemize}

\begin{tcolorbox}[breakable,colback=teal!4,colframe=teal!35!black,title={Before you move on}]
What does \zh{站长} mean? (\textquotedblleft{}Stationmaster; head of a station\textquotedblright{}.) Say it once more.
\end{tcolorbox}
